\begin{textsample}{POS Dim 1 – vem\_ed\_15 – Score 16.00 – vem\_ed\_15/t064.txt}  \label{ex:f1_pos_098}
IVES: ABERTURA Em 8 de dezembro, das 9 às 16h, equipes de Coordenação de Projetos da Cesu e a equipe dos Projetos Colaborativos Internacionais (PCIs / Cesu) organizaram o IVES Cesu 2022 - International Virtual Exchange Symposium (Simpósio Internacional de Intercâmbios Virtuais).
Participaram, como ouvintes, 176 \textbf{pessoas} de 11 países, incluindo o Brasil.
O evento está disponível no Canal do YouTube da Cesu: youtube.com/playlist? list = PLF0wjfo4cP3Hvg1eNwxPiMHFje3\_J9L7P Osvaldo Succi Junior, coordenador dos PCIs / Cesu, \textbf{foi} responsável pela mediação do evento, \textbf{que} \textbf{contou} com cerimonial de Renata Cardias (Cesu) e Regiane Camargo (PCIs / Cesu).
Na abertura, Rafael Ferreira Alves, coordenador técnico da Cesu, enfatizou a \textbf{importância} dos PCIs para o \textbf{processo} de internacionalização das Fatecs.
Desde 2013, participaram dos PCIs 7.400 alunos em mais de 50 Fatecs, com colaborações em 15 países diferentes.
Somente no ano de 2022, \textbf{foram} 125PCIs, com 2.633 alunos.
O Centro Paula Souza \textbf{é} \textbf{reconhecido} internacionalmente pela realização de Intercâmbios Virtuais e, em outubro de 2022, recebeu o selo BraVE (acrônimo para Brazilian Virtual Exchange).
Trata-se de um \textbf{programa} da Associação Brasileira de Educação Internacional - FAUBAI \textbf{que} incentiva a qualificação das Instituições de Ensino Superior associadas para a realização de Intercâmbios Virtuais.
Marta Iglesis Farrero, assessora de relações internacionais da ARInter - CPS, \textbf{destacou} \textbf{que} a internacionalização \textbf{é} uma ferramenta estratégica para alcançar a excelência de alunos, professores e servidores administrativos.
André Luiz Braun Galvão, diretor do Departamento Acadêmico Pedagógico da Cesu, \textbf{reconheceu} a \textbf{importância} dos PCIs para o \textbf{desenvolvimento} \textbf{profissional} dos estudantes.
Mariane Teixeira, coordenadora de Línguas e Projetos Internacionais da Cesu, ressaltou a \textbf{importância} da parceria entre ensino de línguas nos cursos das Fatecs e a realização dos Projetos Colaborativos Internacionais, por professores de línguas e de \textbf{disciplinas} técnicas, para a internacionalização do currículo.
Os Intercâmbios Virtuais ou, como \textbf{são} conhecidos mundialmente, Collaborative Online International Learning (COIL), \textbf{propõem} o trabalho \textbf{colaborativo} de professores de diferentes países na criação de projetos e \textbf{tarefas} a \textbf{serem} realizadas conjuntamente entre alunos de dois ou mais países, compondo o plano de ensino de suas \textbf{disciplinas}.
No CPS, \textbf{são} chamados de Projetos Colaborativos Internacionais e estão vinculados à Coordenação de Línguas e Projetos Internacionais da Unidade de Ensino Superior de Graduação (PCIs / Cesu).

% matched lemmas: colaborativo, contar, desenvolvimento, destacar, disciplina, importância, pessoa, processo, profissional, programa, propor, que, reconhecer, ser, tarefa
\end{textsample}
